\section{Implementation Report}
\label{sec:implementation-report}

This chapter provides a technical description of the interpreter's internal
architecture, the use of ANTLR\,4 for parsing, the stabiliser-based quantum
simulation engine, the multiprocessing model, and the testing framework.

% ─────────────────────────────────────────────
\subsection{Project Structure}
\label{sec:impl-project-structure}

The repository is organised as follows:

\begin{codebox}
QLang/
  antlr/
    QLang.g4              ANTLR4 combined grammar
    antlr-4.13.2-complete.jar
  int/                    Interpreter source (Python)
    QLang/                Generated lexer/parser (output of ANTLR)
    context/              Handler modules for AST node types
      control/            for, while, iterate
      expression/         variables, numbers, lists, booleans
      function/           declaration and invocation
      io/                 print, println, debug, input
      math/               arithmetic and logical operators
      operator/           comparisons, increment, cast, reset
      statement/          statement dispatch, break/continue
      time/               wait
      variable/           declaration, assignment
      network.py          send / receive / available
      quantum.py          gate application and measurement
    exception/            39 runtime exception classes
    operations/           arithmetic helpers
    sim/                  Quantum simulator subsystem
      data/               Data structures (gates, graph, request, ...)
      exceptions/         Simulator-specific exceptions
      simulator.py        Stabiliser tableau core
      server.py           QuantumServer (separate process)
      client.py           QuantumClient (per-place proxy)
    run.py                Entry point
    place.py              Place class (execution unit)
    places.py             Place partitioning logic
    network.py            QuantumNetwork (packet pool)
    scope.py              Scope manager (variable scoping)
    variable.py           Variable class
    expression.py         Expression and ValueResolver
    function.py           Function / FunctionParam
    state.py              State (qubit handle)
    obs.py                Obs / ObsRegister
    num.py                Num
    text.py               Text
    imports.py            Import preprocessor
    consts.py             ANTLR context type aliases
    console.py            Per-place console window
    script_errors.py      Error display (coloured box)
    logger.py             Logging subsystem
  stdlib/                 Standard library (.ql files)
  tests/                  Test framework and test suites
  highlighter/            VS Code syntax highlighting extension
  pyproject.toml          Project metadata and dependencies
  run.bat                 Build-and-run script
\end{codebox}

\subsubsection{Python Dependencies}

The project is managed by the \texttt{uv} tool. Dependencies declared in
\texttt{pyproject.toml}:

\begin{center}
\begin{tabular}{|l|l|p{8cm}|}
\hline
\textbf{Package} & \textbf{Version} & \textbf{Role} \\
\hline
\texttt{antlr4-python3-runtime} & $\geq$ 4.13.2 & Runtime for the generated parser \\
\texttt{colorama}               & $\geq$ 0.4.6  & Coloured terminal output \\
\texttt{pygls}                  & $\geq$ 1.3.0  & Language Server Protocol support \\
\hline
\end{tabular}
\end{center}

The required Python version is \textbf{$\geq$ 3.14}.

% ═══════════════════════════════════════════════
\subsection{ANTLR 4 Grammar and Parser Generation}
\label{sec:impl-antlr}

% ─────────────────────────────────────────────
\subsubsection{Grammar File}

The entire language is defined in a single combined grammar file
\texttt{antlr/QLang.g4}. A combined grammar means the lexer and
parser rules coexist in one file.

Parser generation is executed automatically on every interpreter run by the
\texttt{run.bat} script:

\begin{codebox}
java -jar antlr-4.13.2-complete.jar \
     -Dlanguage=Python3 QLang.g4 \
     -o ..\int\QLang\ -visitor
\end{codebox}

The generated files (\texttt{QLangLexer.py}, \texttt{QLangParser.py},
\texttt{QLangVisitor.py}, and token files) are placed in
\texttt{int/QLang/}.

% ─────────────────────────────────────────────
\subsubsection{Top-Level Grammar Rules}

\begin{codebox}
program : topLevelItem* EOF ;

topLevelItem
    : placeDecl       // place declaration
    | functionDecl    // function (goes to global place)
    | statement       // statement (goes to global place)
    ;
\end{codebox}

A program is a sequence of top-level items: place declarations,
function declarations, or statements. Items that are not wrapped in a
\texttt{place} block belong to the implicit \texttt{global} place.

% ─────────────────────────────────────────────
\subsubsection{Operators}

ANTLR\,4 resolves the ambiguity of left-recursive expressions by assigning
higher precedence to alternatives listed earlier in the \texttt{expr} rule.
The effective precedence table (highest first):

\begin{center}
\begin{tabular}{|c|l|l|}
\hline
\textbf{\#} & \textbf{Operators} & \textbf{Label} \\
\hline
1 & \texttt{!x}, \texttt{-x}, \texttt{+x}   & Unary \\
2 & \texttt{++x}, \texttt{--x}               & Pre-increment / decrement \\
3 & \texttt{x++}, \texttt{x--}               & Post-increment / decrement \\
4 & \texttt{x.0} (float cast)                & FloatCastExpr \\
5 & \texttt{= += -= *= /= \%= **= \&= |=}   & Assignment / compound \\
6 & \texttt{**}                               & Exponentiation \\
7 & \texttt{* / \%}                           & Multiplicative \\
8 & \texttt{+ -}                              & Additive \\
9 & \texttt{>>}                               & Binary-from-list \\
10 & \texttt{< > <= >=}                       & Relational \\
11 & \texttt{== !=}                           & Equality \\
12 & \texttt{\&\&}                            & Logical AND \\
13 & \texttt{||}                              & Logical OR \\
14 & \texttt{\$x}                             & Type query \\
15 & \texttt{c ? a : b}                       & Ternary \\
\hline
\end{tabular}
\end{center}

% ─────────────────────────────────────────────
\subsubsection{Lexer Highlights}

\begin{itemize}
    \item \textbf{Identifiers} use the Unicode pattern
          \texttt{[\textbackslash p\{L\}\_][\textbackslash p\{L\}\textbackslash p\{N\}\_]*},
          supporting non-ASCII characters (including Polish letters).
    \item \textbf{Numbers}: decimal, hexadecimal (\texttt{0x\ldots}),
          binary (\texttt{0b\ldots}), and scientific notation
          (\texttt{1.5e10}).
    \item \textbf{Strings}: three quoting styles --- double quotes, single
          quotes, and backticks.
    \item \textbf{Comments}: line (\texttt{//}) and block
          (\texttt{/* \ldots */}).
\end{itemize}

% ─────────────────────────────────────────────
\subsubsection{Custom Error Listener}

The default ANTLR\,4 error listener is replaced by a custom class,
\texttt{QLangErrorListener} (in \texttt{int/run.py}).  It analyses the
parser's context stack (\texttt{\_ctx}) to produce user-friendly error
messages instead of raw ANTLR output.  Over 50 context-specific error
patterns are defined, each mapping a combination of parser context and
offending token to a human-readable title and explanation (e.g.\
\texttt{LeftOpen}, \texttt{BrokenIf}, \texttt{NoVariable},
\texttt{BadGate}).

% ═══════════════════════════════════════════════
\subsection{Interpreter Pipeline}
\label{sec:impl-pipeline}

The full lifecycle of a QLang program is:

\begin{enumerate}
    \item \textbf{Read} the \texttt{.ql} file (UTF-8).
    \item \textbf{Preprocess imports}: the preprocessor in \texttt{imports.py}
          resolves the \texttt{<< "path" >>} and \texttt{<< ""path"" >>}
          directives before lexing.
    \item \textbf{Lex and parse} with the ANTLR\,4-generated lexer and parser.
    \item \textbf{Partition into places}: \texttt{divideIntoPlaces()} in
          \texttt{places.py} walks the top-level items and groups them by
          place name. For each place, the source text of every member is
          stored together with its original line number.
    \item \textbf{Start the quantum server}: a dedicated
          \texttt{multiprocessing.Process} running
          \texttt{QuantumServer.run()}.
    \item \textbf{Create the packet network}: a shared
          \texttt{QuantumNetwork} backed by a
          \texttt{multiprocessing} manager object.
    \item \textbf{Launch each place} as a separate OS process.
    \item \textbf{Wait} for all place processes to terminate.
    \item \textbf{Shut down} the quantum server.
\end{enumerate}

% ─────────────────────────────────────────────
\subsubsection{Re-Parsing Inside Place Processes}

Because ANTLR\,4 parse trees are not serialisable with Python's
\texttt{pickle}, they cannot be shared across \texttt{multiprocessing}
boundaries.  The interpreter works around this by storing each place's code
as plain text and \textbf{re-parsing} it inside the child process:

\begin{codebox}
# inside Place.process_target()
place_text = "place " + self.name + "{\n"
for block in self.blocks:
    place_text += block.text + "\n"
place_text += "};"

input_stream = InputStream(place_text)
lexer  = QLangLexer(input_stream)
parser = QLangParser(CommonTokenStream(lexer))
tree   = parser.placeDecl()

for member in tree.placeMember():
    self.handle_block(member.getChild(0), None)
\end{codebox}

% ═══════════════════════════════════════════════
\subsection{Handler Dispatch Mechanism}
\label{sec:impl-dispatch}

Although ANTLR\,4 generates a Visitor base class, the interpreter does
not use it.  Instead, a type-to-handler dictionary
is defined in \texttt{Place.handle\_block()}:

\begin{codebox}
HANDLERS = {
    TerminalCtx:          handle_terminal,
    StatementCtx:         handle_statement,
    VarDeclCtx:           handle_variable_declaration,
    GateStmtCtx:          handle_gate_statement,
    MeasureExprCtx:       handle_measure_expr,
    IfStmtCtx:            handle_if,
    ForStmtCtx:           handle_for,
    WhileStmtCtx:         handle_while,
    ...                   # approximately 70 mappings total
}
\end{codebox}

For every AST node, the method iterates over the keys of \texttt{HANDLERS}
and checks \texttt{isinstance(block, }\texttt{type)}.  When a match is found, the
corresponding handler function is called with the signature
\texttt{handler(self, block, parent, pos)}.

Each handler is imported from a dedicated module inside the
\texttt{int/context/} package.  Type aliases for ANTLR context classes
(e.g.\ \texttt{VarDeclCtx = QLangParser.VarDeclContext}) are centralised in
\texttt{int/consts.py}.

% ═══════════════════════════════════════════════
\subsection{Internal Type System}
\label{sec:impl-types}

The interpreter distinguishes two layers of types:

\begin{center}
\begin{tabular}{|l|l|p{7cm}|}
\hline
\textbf{Layer} & \textbf{Types} & \textbf{Purpose} \\
\hline
Primitive &
    \texttt{int}, \texttt{float}, \texttt{bool}, \texttt{string}, \texttt{list} &
    Internal evaluation types carried by \texttt{Expression} \\
\hline
Variable wrapper &
    \texttt{Obs}, \texttt{Num}, \texttt{State}, \texttt{Text}, \texttt{any} &
    Language-level types, each backed by a dedicated Python class \\
\hline
\end{tabular}
\end{center}

\subsubsection{Expression}

\texttt{Expression} (\texttt{int/expression.py}) is the main value carrier
returned by handlers.  It stores a \texttt{type} string, a \texttt{value},
and an optional \texttt{shape}.  The class overloads all Python
arithmetic and comparison dunder methods with automatic type promotion
following the chain: \texttt{bool} $\rightarrow$ \texttt{int} $\rightarrow$
\texttt{float} $\rightarrow$ \texttt{string}.

\subsubsection{ValueResolver}

\texttt{ValueResolver} provides static helper methods:
\begin{itemize}
    \item \texttt{extract\_raw\_value(v)} --- iteratively unwraps nested
          \texttt{Expression}/\texttt{Variable}/wrapper objects down to a Python primitive.
    \item \texttt{resolve\_for\_operation(v)} --- unwraps to the level needed
          for arithmetic, preserving \texttt{Obs} / \texttt{Num} / \texttt{Text}
          wrappers that define their own operators.
    \item \texttt{wrap\_value(type, v)} / \texttt{wrap\_list(\ldots)} ---
          wraps raw values back into the appropriate class.
\end{itemize}

\subsubsection{Variable}

\texttt{Variable} (\texttt{int/variable.py}) supports:
multi-dimensional arrays, dynamic arrays,
constants, three indexing modes
(\texttt{SimpleIndex}, \texttt{RangeIndex}, \texttt{ListIndex}),
character-level indexing into \texttt{Text}, and reset-to-initial-value.

\subsubsection{Scope Manager}

\texttt{ScopeManager} (\texttt{int/scope.py}) implements a
linked-list stack of scopes.  Each scope stores a
\texttt{vars} dictionary, a code position, a type label (\texttt{block},
\texttt{for}, \texttt{if}, \texttt{function}, \ldots), and a parent
reference.  Variable look-up walks from the current scope upward to the
global scope.

% ═══════════════════════════════════════════════
\subsection{Quantum Simulator}
\label{sec:impl-simulator}

\subsubsection{Client--Server Architecture}

The simulator runs in a dedicated OS process.  Communication
between places and the simulator uses \texttt{multiprocessing.Queue}:

\begin{itemize}
    \item One shared command queue --- all places enqueue
          \texttt{QuantumRequest} objects.
    \item One response queue per place --- the server returns
          \texttt{QuantumResponse} to the correct caller.
\end{itemize}

The server runs
an infinite loop that dequeues requests, dispatches them to the
\texttt{QuantumSimulator} instance, and enqueues the response.
Ten command types are supported: \texttt{CREATE}, \texttt{LIST},
\texttt{STAB}, \texttt{STATES}, \texttt{HISTORY}, \texttt{GATE},
\texttt{MEASURE}, \texttt{REMOVE}, \texttt{ALIAS}, \texttt{SEED}.

\subsubsection{Stabiliser Tableau}

The core simulation engine (\texttt{int/sim/simulator.py}) implements the
stabiliser tableau formalism.  For $n$~qubits, the data structure
consists of:

\begin{itemize}
    \item $\mathbf{X}$: a $2n \times n$ bit matrix (X components of Pauli
          operators).
    \item $\mathbf{Z}$: a $2n \times n$ bit matrix (Z components).
    \item $\mathbf{phase}$: a $2n$-element vector (values 0 or 2,
          representing signs $+1$ and $-1$).
\end{itemize}

Rows $0 \ldots n{-}1$ are destabilisers, rows $n \ldots 2n{-}1$
are \emph{stabilisers}.  Each Pauli operator is encoded as a pair of
bits $(x,\,z)$:

\begin{center}
\begin{tabular}{|c|c|c|}
\hline
\textbf{Pauli} & $x$ & $z$ \\
\hline
$I$ & 0 & 0 \\
$X$ & 1 & 0 \\
$Y$ & 1 & 1 \\
$Z$ & 0 & 1 \\
\hline
\end{tabular}
\end{center}

A precomputed Pauli multiplication table stores the product and phase offset for
every pair of Pauli operators, enabling $O(1)$ row-multiplication steps.

\subsubsection{Gate Implementation}

All supported gates belong to the Clifford group and are
implemented as direct bit manipulations on the tableau:

\begin{center}
\begin{tabular}{|l|p{9cm}|}
\hline
\textbf{Gate} & \textbf{Tableau transformation} \\
\hline
$H$ & Swap columns $X \leftrightarrow Z$; adjust phase if both bits are 1. \\
$S$ & $Z_{r,c} \mathrel{\oplus}= X_{r,c}$; adjust phase. \\
$X$ & Flip phase if $Z_{r,c} = 1$. \\
$Y$ & Flip phase if $X_{r,c} \oplus Z_{r,c} = 1$. \\
$Z$ & Flip phase if $X_{r,c} = 1$. \\
\hline
$\mathit{CNOT}(c,t)$ & $X_{r,t} \mathrel{\oplus}= X_{r,c}$;\;
                        $Z_{r,c} \mathrel{\oplus}= Z_{r,t}$;\; phase correction. \\
$\mathit{CZ}$ & Decomposed as $H(t) \to \mathit{CNOT}(c,t) \to H(t)$. \\
\hline
\end{tabular}
\end{center}

The \texttt{swap} gate is decomposed at interpreter level as three consecutive CNOT gates.

Because the stabiliser formalism only supports Clifford gates,
non-Clifford gates (e.g.\ $T$, Toffoli, arbitrary rotations)
cannot be simulated by this engine.  This is a deliberate
trade-off: Clifford circuits can be simulated in $O(n^2)$ time per gate,
whereas a general state-vector simulation would require $O(2^n)$ memory.

\subsubsection{Measurement}

Measurement of qubit $q$ proceeds as follows:

\begin{enumerate}
    \item Search stabiliser rows ($n \ldots 2n{-}1$) for a row where
          $X_{r,q} = 1$.
    \item If found:
    \begin{itemize}
        \item The outcome is sampled uniformly ($0$ or $1$).  If a per-prefix
              seed was set via \texttt{seed()}, a deterministic
              \texttt{random.Random(seed)} is used.
        \item The tableau is updated to reflect the post-measurement
              projective collapse.
    \end{itemize}
    \item If not found:
    \begin{itemize}
        \item The outcome is computed from the destabiliser rows.
    \end{itemize}
    \item The entanglement graph calls \texttt{split(q)} to remove $q$ from
          its entanglement group.
\end{enumerate}

\subsubsection{Entanglement Graph}

The class \texttt{EntaglementGraph} (\texttt{int/sim/data/graph.py})
implements a \textbf{Union--Find} structure with path compression and
union-by-rank.  It tracks which qubits are entangled:

\begin{itemize}
    \item \texttt{union(a, b)} --- called when a two-qubit gate links
          $a$ and $b$.
    \item \texttt{split(x)} --- called after measurement, disconnects
          $x$ from its group.
    \item \texttt{connected(a, b)} --- checks whether two qubits share an
          entanglement group.
    \item \texttt{remove(x)} --- used by \texttt{remove\_qubit}, returns
          \texttt{False} if $x$ is still entangled.
\end{itemize}

\subsubsection{State Aliasing}

The method \texttt{rename\_state(old\_id, new\_prefix)} creates an
alias for an existing simulator state under a new prefix.  This is essential
when qubits are transferred between places: the received qubit gets a new
alias with the receiver's prefix but still references the same simulator
state.

% ═══════════════════════════════════════════════
\subsection{Multiprocessing Model}
\label{sec:impl-multiprocessing}

Each place runs as a separate OS process
(\texttt{multiprocessing.Process}).

Key consequences:
\begin{itemize}
    \item Parse trees must be re-created in each process.
    \item Shared state requires IPC primitives like
          \texttt{multiprocessing} \texttt{.Queue} and
          \texttt{multiprocessing} \texttt{.Manager}.
\end{itemize}

\subsubsection{Packet Network Internals}

The \texttt{QuantumNetwork} class (\texttt{int/network.py}) is backed by:

\begin{codebox}
manager      = multiprocessing.Manager()
packet_pool  = manager.list()                    # shared packet list
condition    = manager.Condition(manager.Lock())  # for wait/notify
global_counter = manager.Value('i', 0)           # monotonic packet ID
\end{codebox}

Sending (\texttt{send}) acquires the condition lock, appends a
\texttt{Packet} to \texttt{packet\_pool}, increments the counter, and calls
\texttt{condition.notify\_all()}.

Receiving (\texttt{wait\_for}) acquires the condition lock and
scans the pool for a matching packet.  If no match is found, it calls
\texttt{condition.wait()} and re-scans when notified.  Matching considers:
quantum/classical flag, data size, target ID, source ID, and message name.

Peeking (\texttt{peek}) performs the same scan but does not consume
or block.

\subsubsection{Per-Place Console}

Each place opens a separate CMD window via the \texttt{Console} class
(\texttt{int/console.py}).  The window is a subprocess running
\texttt{console\_worker.py}, connected to the place process through a
TCP socket on \texttt{localhost}.

% ═══════════════════════════════════════════════
\subsection{Error Reporting Internals}
\label{sec:impl-errors}

The \texttt{ScriptErrors.showError()} method (\texttt{int/script\_errors.py})
renders a coloured, box-drawn error display.

\subsubsection{Runtime Error Verification}

For tests that expect a specific runtime error, the framework parses
structured \texttt{ERROR\_*} entries from the log:

\begin{codebox}
[--- TEST ---](D) ERROR_LINE_START[0]=5
[--- TEST ---](D) ERROR_LINE_END[0]=5
[--- TEST ---](D) ERROR_COL_START[0]=10
[--- TEST ---](D) ERROR_COL_END[0]=15
[--- TEST ---](D) ERROR_CODE[0]=DBZ-1
\end{codebox}

The helper \texttt{find\_error\_entry(line\_start,} \texttt{line\_end,}
\texttt{col\_start,} \texttt{col\_end,} \texttt{code)} matches against these entries, allowing tests to verify
both the error code and the exact source location.

\subsubsection{Syntax Error Collection Mode}

A special mode (\texttt{SYNTAX\_ERROR\_COLLECT}) scans all \texttt{.ql}
files in \texttt{tests/syntax\_errors/}, runs the interpreter in
\texttt{SYNTAX\_MODE} (parse-only, no execution), and collects error data
into \texttt{tests/} \texttt{syntax\_errors.json}.  This is used to validate that
the custom \texttt{QLangErrorListener} produces meaningful messages for
every syntax error pattern.

% ═══════════════════════════════════════════════
\subsection{Key Architectural Decisions}
\label{sec:impl-decisions}

\begin{center}
\begin{tabular}{|p{4.5cm}|p{8.5cm}|}
\hline
\textbf{Decision} & \textbf{Rationale} \\
\hline
Stabiliser tableau instead of state vector &
    $O(n^2)$ per-gate cost vs.\ $O(2^n)$ memory for state vectors;
    sufficient for the supported Clifford gate set. \\
\hline
OS processes instead of threads &
    Bypasses Python's GIL, giving genuine parallelism for places. \\
\hline
Re-parsing per place &
    ANTLR\,4 CST objects are not picklable; storing text and re-parsing
    is the simplest workaround. \\
\hline
Handler dictionary instead of ANTLR Visitor &
    Enables modular, file-per-handler organisation; avoids a monolithic
    Visitor class. \\
\hline
TCP sockets for console I/O &
    Connects the place process to a separate CMD window process on
    Windows, allows independent lifecycle management. \\
\hline
\end{tabular}
\end{center}

% ═══════════════════════════════════════════════
\subsection{Limitations}
\label{sec:impl-limitations}

\begin{itemize}
    \item \textbf{No non-Clifford gates.}  The stabiliser formalism cannot
          simulate gates outside the Clifford group (e.g.\ $T$, Toffoli,
          arbitrary rotations).  Universal quantum computation is therefore
          not expressible.
    \item \textbf{Re-parsing overhead.}  Each place re-parses its code at
          start-up, which adds latency proportional to the code size.
\end{itemize}
